Model zoos contain thousands of trained neural networks, but methods for predicting \emph{how} models behave---beyond scalar accuracy---remain underdeveloped. We investigate whether class-wise behavioral profiles can be extracted from weight matrices alone. On the Small CNN Zoo (193 CIFAR-10 models), we find that 67.6\% of class-wise accuracy variance is unexplained by overall accuracy---exceeding our detection threshold by 13.5$\times$. This establishes that behavioral fingerprints exist: models with identical accuracy fail on different classes in distinct patterns. However, simple weight statistics (per-layer mean, std, norm) fail to extract this signal, achieving $R^2 = -0.08$ compared to a stratified baseline $R^2 = 0.12$. Our negative result on the mechanism ($\Delta R^2 = -0.20$) combined with the positive existence finding motivates learned representations---NF-Layers, SANE embeddings, or behavioral autoencoders---as the path toward behavioral prediction from weights.
